% 厦门大学博士后研究工作报告版式
\RequirePackage{geometry}
\RequirePackage{setspace}
\RequirePackage{tabularx,booktabs}
\RequirePackage{graphicx}
\RequirePackage{amsmath,amssymb}
\RequirePackage{caption}
\RequirePackage[titles]{tocloft}
\RequirePackage{fancyhdr}
\RequirePackage[hidelinks,bookmarksnumbered]{hyperref}
\RequirePackage[backend=biber,style=gb7714-2015,sorting=none]{biblatex}

% A4 版心采用博士后报告编写规则，字号遵循厦大范本。
\geometry{left=35mm,right=20mm,top=30mm,bottom=25mm,footskip=12mm}
\IfFontExistsTF{Times New Roman}
  {\setmainfont{Times New Roman}}
  {\setmainfont{texgyretermes}[
    Extension=.otf,
    UprightFont=*-regular,
    BoldFont=*-bold,
    ItalicFont=*-italic,
    BoldItalicFont=*-bolditalic
  ]}
\IfFontExistsTF{SimSun}
  {\setCJKmainfont[Script=CustomDefault,AutoFakeBold,AutoFakeSlant]{SimSun}}
  {\IfFontExistsTF{Songti SC}
    {\setCJKmainfont[Script=CustomDefault,AutoFakeSlant]{Songti SC}}
    {\setCJKmainfont[Script=CustomDefault,BoldFont=FandolSong-Bold.otf,ItalicFont=FandolKai-Regular.otf]{FandolSong-Regular.otf}}}
\IfFontExistsTF{SimHei}
  {\setCJKsansfont[Script=CustomDefault,AutoFakeBold]{SimHei}}
  {\setCJKsansfont[Script=CustomDefault,BoldFont=FandolHei-Bold.otf]{FandolHei-Regular.otf}}

\AtBeginDocument{\setstretch{1.5}\zihao{-4}}
\setlength{\parindent}{2\ccwd}
\setlength{\parskip}{0pt}

\ctexset{
  contentsname = {目\quad 次},
  chapter = {
    name = {,}, number = \arabic{chapter},
    format = \centering\sffamily\bfseries\zihao{3},
    aftername = \quad,
    beforeskip = 0pt, afterskip = 18pt,
    afterindent = true, pagestyle = empty
  },
  section = {
    format = \sffamily\bfseries\zihao{4},
    beforeskip = 12pt, afterskip = 6pt,
    afterindent = true
  },
  subsection = {
    format = \sffamily\bfseries\zihao{-4},
    beforeskip = 8pt, afterskip = 4pt,
    afterindent = true
  },
  appendix = {name = {附录\space,}, number = \Alph{chapter}}
}
\setcounter{tocdepth}{2}
\setcounter{secnumdepth}{2}
\cftsetindents{chapter}{0pt}{3.5em}
\cftsetindents{section}{1em}{2.5em}
\cftsetindents{subsection}{2em}{3.2em}
\cftsetpnumwidth{2.8em}
\cftsetrmarg{3.5em}
\setlength{\cftbeforechapskip}{0pt}
\renewcommand{\cftchapfont}{\normalfont}
\renewcommand{\cftchappagefont}{\normalfont}
\renewcommand{\cftchapleader}{\cftdotfill{\cftdotsep}}

\fancypagestyle{plain}{
  \fancyhf{}
  \fancyfoot[R]{\zihao{-4}\thepage}
  \renewcommand{\headrulewidth}{0pt}
}
\newcommand{\XmuPostdocMainMatter}{%
  \clearpage
  \pagenumbering{arabic}%
  \hypersetup{pageanchor=true}%
  \ctexset{chapter/pagestyle=plain}%
  \pagestyle{plain}%
}

\captionsetup{font=small,labelsep=space}
\captionsetup[figure]{position=bottom}
\captionsetup[table]{position=top}
\numberwithin{equation}{chapter}
\newcommand{\XmuPostdocAppendix}{%
  \appendix
  \renewcommand{\thefigure}{\Alph{chapter}\arabic{figure}}%
  \renewcommand{\thetable}{\Alph{chapter}\arabic{table}}%
  \renewcommand{\theequation}{\Alph{chapter}\arabic{equation}}%
}
\newcommand{\XmuPostdocBackChapter}[1]{%
  \chapter*{#1}%
  \addcontentsline{toc}{chapter}{#1}%
}

% 封面和题名页采用对称版心。标题和较长的元信息可自动折行。
\newenvironment{XmuPostdocTitleSheet}{%
  \newgeometry{left=28mm,right=28mm,top=30mm,bottom=29mm}%
  \thispagestyle{empty}%
  \centering
}{%
  \par\restoregeometry
}
\newcommand{\XmuPostdocTitleField}[2]{#1 & #2\\}
\newcommand{\XmuPostdocBlank}[1]{\rule[-.3ex]{#1}{.4pt}}
